\documentclass[11pt]{article}
\usepackage[margin=1in]{geometry}
\usepackage{enumitem}
\setlist{nosep}

\title{Live AR Sudoku Solver}
\author{}
\date{}

\begin{document}
\maketitle

\textbf{Goal.} Point a webcam at a printed Sudoku. The app detects the grid, reads the digits,
solves the puzzle and overlays the solution on the paper in real time, following the paper
as it moves.

\textbf{Pipeline.}
\begin{enumerate}
  \item Grid detection: adaptive threshold, Canny edges, largest 4-sided contour, and Hough
        lines to refine the grid.
  \item Rectification: homography from the 4 corners to a flat square image.
  \item Cell reading: split into 81 cells, detect empty cells, classify digits with a small
        CNN trained on MNIST plus printed-font digits.
  \item Solving: backtracking with constraint propagation.
  \item AR overlay: draw the solution through the inverse homography and track the corners
        across frames instead of re-detecting every frame.
\end{enumerate}

\textbf{Beyond existing hobby versions.} Simple versions exist online, so we focus on:
\begin{itemize}
  \item robustness to tilt, glare, shadows and curved paper;
  \item a stable overlay from frame to frame;
  \item a measured comparison of contour-based vs Hough-based grid detection.
\end{itemize}

\textbf{Evaluation.} About 50 photos and 10 videos under varied conditions. Metrics: grid
detection rate, digit accuracy per cell, fully-solved rate, FPS, overlay jitter.

\textbf{Course fit.} Edges, Hough transform, thresholding and segmentation, homography,
recognition.

\textbf{Risk.} Low: every stage has a known baseline, so effort goes into robustness and the
demo.

\end{document}
